% GENERATED FILE. Edit canonical lessons, then run npm run generate:books.
% canonical-source-hash: fnv1a64:8daadbb27fcfb629
% canonical-lessons: TA-C399-eppatiyum, TA-C399-kututal, TA-C399-mulam, TA-C399-ullatakku, TA-C399-kariyam, TA-R399-review-393-399, TA-R399-review-final-399

\chapter{Whatever Happens}
\label{ch:ta-whatever-happens}
\hlchaptermodality{399}

\begin{chapteropening}
\textbf{By the end of this chapter:} \emph{I can say that something will happen anyway, through someone, or include something.}

Complete the last lesson of chapter 399: I can say that something will happen anyway, through someone, or include something.
\end{chapteropening}

\section[\texorpdfstring{\ta{எப்படியும்} (eppaṭiyum)}{எப்படியும் (eppaṭiyum)}]{\ta{எப்படியும்} (eppaṭiyum) — anyway; in any case}

\label{lesson:TA-C399-eppatiyum}



\begin{quote}
Before the new one: say the Tamil for stability, then the Tamil for a matter.
\end{quote}

\subsection*{You'll want to know: \ta{எப்படியும்}}
\textbf{\ta{எப்படியும்}} — \emph{eppaṭiyum} — \textquotedblleft{}anyway; in any case\textquotedblright{}.

\begin{quote}
\ta{மழை} \ta{பெய்தாலும்} \ta{எப்படியும்} \ta{நான்} \ta{வருவேன்}.
\end{quote}

\begin{quote}
\emph{maḻai peytālum eppaṭiyum nāṉ varuvēṉ.}
\end{quote}

\begin{quote}
Even if it rains, I will come anyway.
\end{quote}

\ta{எப்படி}, how + -\ta{உம்}: however it goes.

\begin{cousinweb}
\textbf{\ta{எப்படியும்}} (\emph{eppaṭiyum}) is native: \emph{eppaṭi}, \textquotedblleft{}how\textquotedblright{}, built from the question root \emph{e-} (DEDR 5151) and \emph{paṭi}, \textquotedblleft{}manner\textquotedblright{} (DEDR 3851), with the particle -\ta{உம்}, \textquotedblleft{}even\textquotedblright{}.
\end{cousinweb}

\begin{grammarlens}[title={What you've built}]
One more word for saying what happens if.
\end{grammarlens}

\subsection*{Your turn}
\begin{itemize}
\raggedright
  \item \emph{Say it:} \emph{eppaṭiyum}
  \item \emph{Say it:} \emph{eppaṭiyum}, once more
  \item \emph{Say it:} \emph{maḻai peytālum eppaṭiyum nāṉ varuvēṉ.} — Even if it rains, I will come anyway.
  \item \emph{Recall:} say the Tamil for stability, then the Tamil for a matter, then say \emph{eppaṭiyum} again
\end{itemize}

\begin{tcolorbox}[breakable,colback=teal!4,colframe=teal!35!black,title={Before you move on}]
What does \ta{எப்படியும்} mean? (\textquotedblleft{}Anyway; in any case\textquotedblright{}.) Say it once more.
\end{tcolorbox}
\section[\texorpdfstring{\ta{கூடுதல்} (kūṭutal)}{கூடுதல் (kūṭutal)}]{\ta{கூடுதல்} (kūṭutal) — additional; extra}

\label{lesson:TA-C399-kututal}



\begin{quote}
Before the new one: say the Tamil for a matter, then the Tamil for anyway.
\end{quote}

\subsection*{You'll want to know: \ta{கூடுதல்}}
\textbf{\ta{கூடுதல்}} — \emph{kūṭutal} — \textquotedblleft{}additional; extra\textquotedblright{}.

\begin{quote}
\ta{கூடுதல்} \ta{பணம்} \ta{கொடுத்தால்} \ta{சீக்கிரம்} \ta{கிடைக்கும்}.
\end{quote}

\begin{quote}
\emph{kūṭutal paṇam koṭuttāl cīkkiram kiṭaikkum.}
\end{quote}

\begin{quote}
If you pay extra, you will get it sooner.
\end{quote}

Before a noun: \ta{கூடுதல்} \ta{கட்டணம்}, an extra charge; \ta{கூடுதலாக}, in addition.

\begin{cousinweb}
\textbf{\ta{கூடுதல்}} (\emph{kūṭutal}) is native Dravidian: DEDR 1882 lists Tamil \emph{kūṭu}, \textquotedblleft{}to come together, be stored up\textquotedblright{}, and \emph{kūṭutal}, \textquotedblleft{}increase\textquotedblright{}.
\end{cousinweb}

\begin{grammarlens}[title={What you've built}]
One more word for saying what happens if.
\end{grammarlens}

\subsection*{Your turn}
\begin{itemize}
\raggedright
  \item \emph{Say it:} \emph{kūṭutal}
  \item \emph{Say it:} \emph{kūṭutal}, once more
  \item \emph{Say it:} \emph{kūṭutal paṇam koṭuttāl cīkkiram kiṭaikkum.} — If you pay extra, you will get it sooner.
  \item \emph{Recall:} say the Tamil for a matter, then the Tamil for anyway, then say \emph{kūṭutal} again
\end{itemize}

\begin{tcolorbox}[breakable,colback=teal!4,colframe=teal!35!black,title={Before you move on}]
What does \ta{கூடுதல்} mean? (\textquotedblleft{}Additional; extra\textquotedblright{}.) Say it once more.
\end{tcolorbox}
\section[\texorpdfstring{\ta{மூலம்} (mūlam)}{மூலம் (mūlam)}]{\ta{மூலம்} (mūlam) — through; by means of; a source}

\label{lesson:TA-C399-mulam}



\begin{quote}
Before the new one: say the Tamil for anyway, then the Tamil for additional.
\end{quote}

\subsection*{You'll want to know: \ta{மூலம்}}
\textbf{\ta{மூலம்}} — \emph{mūlam} — \textquotedblleft{}through; by means of; a source\textquotedblright{}.

\begin{quote}
\ta{இந்தச்} \ta{செய்தி} \ta{எனக்கு} \ta{ஒரு} \ta{நண்பர்} \ta{மூலம்} \ta{தெரிந்தது}.
\end{quote}

\begin{quote}
\emph{intac ceyti eṉakku oru naṇpar mūlam terintatu.}
\end{quote}

\begin{quote}
I learned this news through a friend.
\end{quote}

After a noun: X \ta{மூலம்}, through X. As a noun, a source: \ta{நதியின்} \ta{மூலம்}, the source of a river.

\begin{cousinweb}
\textbf{\ta{மூலம்}} (\emph{mūlam}) is borrowed from Sanskrit \emph{mūla}, \textquotedblleft{}root\textquotedblright{} (CDIAL 10250): the root a thing comes through.
\end{cousinweb}

\begin{grammarlens}[title={What you've built}]
One more word for saying what happens if.
\end{grammarlens}

\subsection*{Your turn}
\begin{itemize}
\raggedright
  \item \emph{Say it:} \emph{mūlam}
  \item \emph{Say it:} \emph{mūlam}, once more
  \item \emph{Say it:} \emph{intac ceyti eṉakku oru naṇpar mūlam terintatu.} — I learned this news through a friend.
  \item \emph{Recall:} say the Tamil for anyway, then the Tamil for additional, then say \emph{mūlam} again
\end{itemize}

\begin{tcolorbox}[breakable,colback=teal!4,colframe=teal!35!black,title={Before you move on}]
What does \ta{மூலம்} mean? (\textquotedblleft{}Through; by means of; a source\textquotedblright{}.) Say it once more.
\end{tcolorbox}
\section[\texorpdfstring{\ta{உள்ளடக்கு} (uḷḷaṭakku)}{உள்ளடக்கு (uḷḷaṭakku)}]{\ta{உள்ளடக்கு} (uḷḷaṭakku) — to include; to contain}

\label{lesson:TA-C399-ullatakku}



\begin{quote}
Before the new one: say the Tamil for additional, then the Tamil for through.
\end{quote}

\subsection*{You'll want to know: \ta{உள்ளடக்கு}}
\textbf{\ta{உள்ளடக்கு}} — \emph{uḷḷaṭakku} — \textquotedblleft{}to include; to contain\textquotedblright{}.

\begin{quote}
\ta{இந்த} \ta{விலை} \ta{உணவையும்} \ta{உள்ளடக்கியது}.
\end{quote}

\begin{quote}
\emph{inta vilai uṇavaiyum uḷḷaṭakkiyatu.}
\end{quote}

\begin{quote}
This price includes food.
\end{quote}

Past (it …): \ta{உள்ளடக்கியது}, \emph{uḷḷaṭakkiyatu}.

Mostly in this form: X-\ta{ஐ} \ta{உள்ளடக்கியது}, it includes X. \ta{உள்ளடக்கம்}, contents.

\begin{cousinweb}
\textbf{\ta{உள்ளடக்கு}} (\emph{uḷḷaṭakku}) is native: \emph{uḷ}, \textquotedblleft{}inside\textquotedblright{} (DEDR 698), and \emph{aṭakku}, \textquotedblleft{}to control, hold in, contain\textquotedblright{} (DEDR 63): to hold inside.
\end{cousinweb}

\begin{grammarlens}[title={What you've built}]
One more word for saying what happens if.
\end{grammarlens}

\subsection*{Your turn}
\begin{itemize}
\raggedright
  \item \emph{Say it:} \emph{uḷḷaṭakku}
  \item \emph{Say it:} \emph{uḷḷaṭakku}, once more
  \item \emph{Say it:} \emph{inta vilai uṇavaiyum uḷḷaṭakkiyatu.} — This price includes food.
  \item \emph{Recall:} say the Tamil for additional, then the Tamil for through, then say \emph{uḷḷaṭakku} again
\end{itemize}

\begin{tcolorbox}[breakable,colback=teal!4,colframe=teal!35!black,title={Before you move on}]
What does \ta{உள்ளடக்கு} mean? (\textquotedblleft{}To include; to contain\textquotedblright{}.) Say it once more.
\end{tcolorbox}
\section[\texorpdfstring{\ta{காரியம்} (kāriyam)}{காரியம் (kāriyam)}]{\ta{காரியம்} (kāriyam) — a task; an affair}

\label{lesson:TA-C399-kariyam}



\begin{quote}
Before the new one: say the Tamil for through, then the Tamil for to include.
\end{quote}

\subsection*{You'll want to know: \ta{காரியம்}}
\textbf{\ta{காரியம்}} — \emph{kāriyam} — \textquotedblleft{}a task; an affair\textquotedblright{}.

\begin{quote}
\ta{இந்தக்} \ta{காரியம்} \ta{முடிந்தால்} \ta{நாம்} \ta{வீட்டுக்குப்} \ta{போகலாம்}.
\end{quote}

\begin{quote}
\emph{intak kāriyam muṭintāl nām vīṭṭukkup pōkalām.}
\end{quote}

\begin{quote}
If this task is done, we can go home.
\end{quote}

\ta{காரியம்} \ta{ஆனது}, it worked, the job got done.

\begin{cousinweb}
\textbf{\ta{காரியம்}} (\emph{kāriyam}) is borrowed from Sanskrit \emph{kārya}, \textquotedblleft{}to be done; business\textquotedblright{} (CDIAL 3078). Tamil breaks the cluster ry with a vowel, as Pali \emph{kāriya} does.
\end{cousinweb}

\begin{grammarlens}[title={What you've built}]
One more word for saying what happens if.
\end{grammarlens}

\subsection*{Your turn}
\begin{itemize}
\raggedright
  \item \emph{Say it:} \emph{kāriyam}
  \item \emph{Say it:} \emph{kāriyam}, once more
  \item \emph{Say it:} \emph{intak kāriyam muṭintāl nām vīṭṭukkup pōkalām.} — If this task is done, we can go home.
  \item \emph{Recall:} say the Tamil for through, then the Tamil for to include, then say \emph{kāriyam} again
\end{itemize}

\begin{tcolorbox}[breakable,colback=teal!4,colframe=teal!35!black,title={Before you move on}]
What does \ta{காரியம்} mean? (\textquotedblleft{}A task; an affair\textquotedblright{}.) Say it once more.
\end{tcolorbox}
\section[(dialogue)]{Looking back: conditions and possibilities}

\label{lesson:TA-R399-review-393-399}



\begin{quote}
Say the Tamil for a report and a census.
\end{quote}

\subsection*{Your turn}
Say these aloud:

\begin{itemize}
\raggedright
  \item a report, a census, to find out, roughly, mostly
  \item one-sidedness, a hidden motive, neutrality, on behalf of, about
  \item to conduct, to create, to remove, to abolish, to abandon
  \item to incite, a reaction, to impose, to reflect, to take into account
  \item democracy, a public debate, unanimously, personal, rational thought
  \item a trend, a challenge, an achievement, stability, a matter
  \item anyway, additional, through, to include, a task
\end{itemize}

\begin{tcolorbox}[breakable,colback=teal!4,colframe=teal!35!black,title={Before you move on}]
A report? (\textbf{\ta{அறிக்கை}}, \emph{aṟikkai}.) A census? (\textbf{\ta{கணக்கெடுப்பு}}, \emph{kaṇakkeṭuppu}.) To find out? (\textbf{\ta{கண்டறி}}, \emph{kaṇṭaṟi}.) Roughly? (\textbf{\ta{சுமார்}}, \emph{cumār}.) Mostly? (\textbf{\ta{பெரும்பாலும்}}, \emph{perumpālum}.) One-sidedness? (\textbf{\ta{ஒருதலைப்பட்சம்}}, \emph{orutalaippaṭcam}.) A hidden motive? (\textbf{\ta{உள்நோக்கம்}}, \emph{uḷnōkkam}.) Neutrality? (\textbf{\ta{நடுநிலை}}, \emph{naṭunilai}.) On behalf of? (\textbf{\ta{சார்பாக}}, \emph{cārpāka}.) About? (\textbf{\ta{குறித்து}}, \emph{kuṟittu}.) To conduct? (\textbf{\ta{நடத்து}}, \emph{naṭattu}.) To create? (\textbf{\ta{உருவாக்கு}}, \emph{uruvākku}.) To remove? (\textbf{\ta{நீக்கு}}, \emph{nīkku}.) To abolish? (\textbf{\ta{ஒழி}}, \emph{oḻi}.) To abandon? (\textbf{\ta{கைவிடு}}, \emph{kaiviṭu}.) To incite? (\textbf{\ta{தூண்டு}}, \emph{tūṇṭu}.) A reaction? (\textbf{\ta{எதிர்வினை}}, \emph{etirviṉai}.) To impose? (\textbf{\ta{சுமத்து}}, \emph{cumattu}.) To reflect? (\textbf{\ta{பிரதிபலி}}, \emph{piratipali}.) To take into account? (\textbf{\ta{பொருட்படுத்து}}, \emph{poruṭpaṭuttu}.) Democracy? (\textbf{\ta{மக்களாட்சி}}, \emph{makkaḷāṭci}.) A public debate? (\textbf{\ta{பட்டிமன்றம்}}, \emph{paṭṭimaṉṟam}.) Unanimously? (\textbf{\ta{ஒருமனதாக}}, \emph{orumaṉatāka}.) Personal? (\textbf{\ta{தனிப்பட்ட}}, \emph{taṉippaṭṭa}.) Rational thought? (\textbf{\ta{பகுத்தறிவு}}, \emph{pakuttaṟivu}.) A trend? (\textbf{\ta{போக்கு}}, \emph{pōkku}.) A challenge? (\textbf{\ta{சவால்}}, \emph{cavāl}.) An achievement? (\textbf{\ta{சாதனை}}, \emph{cātaṉai}.) Stability? (\textbf{\ta{நிலைத்தன்மை}}, \emph{nilaittaṉmai}.) A matter? (\textbf{\ta{விவகாரம்}}, \emph{vivakāram}.) Anyway? (\textbf{\ta{எப்படியும்}}, \emph{eppaṭiyum}.) Additional? (\textbf{\ta{கூடுதல்}}, \emph{kūṭutal}.) Through? (\textbf{\ta{மூலம்}}, \emph{mūlam}.) To include? (\textbf{\ta{உள்ளடக்கு}}, \emph{uḷḷaṭakku}.) A task? (\textbf{\ta{காரியம்}}, \emph{kāriyam}.)
\end{tcolorbox}
\section[(dialogue)]{Whatever Happens — once more}

\label{lesson:TA-R399-review-final-399}



\begin{quote}
Say the Tamil for anyway and additional.
\end{quote}

\subsection*{Your turn}
\begin{itemize}
\raggedright
  \item \emph{Say it:} anyway, additional, through, to include, a task
\end{itemize}

\begin{tcolorbox}[breakable,colback=teal!4,colframe=teal!35!black,title={Before you move on}]
Anyway? (\textbf{\ta{எப்படியும்}}, \emph{eppaṭiyum}.) Additional? (\textbf{\ta{கூடுதல்}}, \emph{kūṭutal}.) Through? (\textbf{\ta{மூலம்}}, \emph{mūlam}.) To include? (\textbf{\ta{உள்ளடக்கு}}, \emph{uḷḷaṭakku}.) A task? (\textbf{\ta{காரியம்}}, \emph{kāriyam}.)
\end{tcolorbox}
